% GENERATED FILE. Edit canonical lessons, then run npm run generate:books.
% canonical-source-hash: fnv1a64:c6c345f5d9c37539
% canonical-lessons: HI-C315-dava, HI-C315-janta, HI-C315-ghosna, HI-C315-jankari, HI-C315-ttaalna

\chapter{Claims and Denials}
\label{ch:hi-claims-and-denials}
\hlchaptermodality{315}

\begin{chapteropening}
\textbf{By the end of this chapter:} \emph{I can report what someone claimed and how it was answered.}

Complete the last lesson of chapter 315: I can report what someone claimed and how it was answered.
\end{chapteropening}

\section[\texorpdfstring{\hi{दावा} (dāvā)}{दावा (dāvā)}]{\hi{दावा} (dāvā) — a claim}

\label{lesson:HI-C315-dava}



\begin{quote}
Before the new one: say the Hindi for to stir up, then the Hindi for a sensation.
\end{quote}

\subsection*{You'll want to know: \hi{दावा}}
\textbf{\hi{दावा}} — \emph{dāvā} — \textquotedblleft{}a claim\textquotedblright{}.

\begin{quote}
\hi{मंत्री} \hi{ने} \hi{दावा} \hi{किया} \hi{कि} \hi{सब} \hi{ठीक} \hi{है।}
\end{quote}

\begin{quote}
\emph{mantrī ne dāvā kiyā ki sab ṭhīk hai.}
\end{quote}

\begin{quote}
The minister claimed that everything was fine.
\end{quote}

Masculine.

\begin{cousinweb}
\textbf{\hi{दावा}} (\emph{dāvā}) is usually traced, through Persian, to Arabic \emph{daʿwā}, \textquotedblleft{}a claim, a lawsuit\textquotedblright{}. \textbf{\hi{दावा} \hi{करना}} is to claim; the claim is masculine: \textbf{\hi{दावा} \hi{किया}}.
\end{cousinweb}

\begin{grammarlens}[title={What you've built}]
One more word for telling what happened.
\end{grammarlens}

\subsection*{Your turn}
\begin{itemize}
\raggedright
  \item \emph{Say it:} \emph{dāvā}
  \item \emph{Say it:} \emph{dāvā}, once more
  \item \emph{Say it:} \emph{mantrī ne dāvā kiyā ki sab ṭhīk hai.} — The minister claimed that everything was fine.
  \item \emph{Recall:} say the Hindi for to stir up, then the Hindi for a sensation, then say \emph{dāvā} again
\end{itemize}

\begin{tcolorbox}[breakable,colback=teal!4,colframe=teal!35!black,title={Before you move on}]
What does \hi{दावा} mean? (\textquotedblleft{}A claim\textquotedblright{}.) Say it once more.
\end{tcolorbox}
\section[\texorpdfstring{\hi{जनता} (jantā)}{जनता (jantā)}]{\hi{जनता} (jantā) — the public; the people}

\label{lesson:HI-C315-janta}



\begin{quote}
Before the new one: say the Hindi for a sensation, then the Hindi for a claim.
\end{quote}

\subsection*{You'll want to know: \hi{जनता}}
\textbf{\hi{जनता}} — \emph{jantā} — \textquotedblleft{}the public; the people\textquotedblright{}.

\begin{quote}
\hi{जनता} \hi{उस} \hi{दावे} \hi{से} \hi{ख़ुश} \hi{नहीं} \hi{थी।}
\end{quote}

\begin{quote}
\emph{jantā us dāve se khuś nahīṃ thī.}
\end{quote}

\begin{quote}
The public was not happy with that claim.
\end{quote}

Feminine.

\begin{cousinweb}
\textbf{\hi{जनता}} (\emph{jantā}) is Sanskrit \emph{janatā}, \textquotedblleft{}a number of people, a community\textquotedblright{} (Monier-Williams), from \emph{jana}, a person. It is singular and feminine in Hindi: \textbf{\hi{जनता} \hi{ख़ुश} \hi{नहीं} \hi{थी}}.
\end{cousinweb}

\begin{grammarlens}[title={What you've built}]
One more word for telling what happened.
\end{grammarlens}

\subsection*{Your turn}
\begin{itemize}
\raggedright
  \item \emph{Say it:} \emph{jantā}
  \item \emph{Say it:} \emph{jantā}, once more
  \item \emph{Say it:} \emph{jantā us dāve se khuś nahīṃ thī.} — The public was not happy with that claim.
  \item \emph{Recall:} say the Hindi for a sensation, then the Hindi for a claim, then say \emph{jantā} again
\end{itemize}

\begin{tcolorbox}[breakable,colback=teal!4,colframe=teal!35!black,title={Before you move on}]
What does \hi{जनता} mean? (\textquotedblleft{}The public; the people\textquotedblright{}.) Say it once more.
\end{tcolorbox}
\section[\texorpdfstring{\hi{घोषणा} (ghoṣṇā)}{घोषणा (ghoṣṇā)}]{\hi{घोषणा} (ghoṣṇā) — a declaration; an official announcement}

\label{lesson:HI-C315-ghosna}



\begin{quote}
Before the new one: say the Hindi for a claim, then the Hindi for the public.
\end{quote}

\subsection*{You'll want to know: \hi{घोषणा}}
\textbf{\hi{घोषणा}} — \emph{ghoṣṇā} — \textquotedblleft{}a declaration; an official announcement\textquotedblright{}.

\begin{quote}
\hi{सरकार} \hi{ने} \hi{नई} \hi{सड़क} \hi{की} \hi{घोषणा} \hi{की।}
\end{quote}

\begin{quote}
\emph{sarkār ne naī saṛak kī ghoṣṇā kī.}
\end{quote}

\begin{quote}
The government announced a new road.
\end{quote}

Feminine.

\begin{cousinweb}
\textbf{\hi{घोषणा}} (\emph{ghoṣṇā}) is Sanskrit \emph{ghoṣaṇā}, \textquotedblleft{}a proclaiming aloud\textquotedblright{} (Monier-Williams), from \emph{ghoṣa}, a loud noise. It is the formal word beside \textbf{\hi{ऐलान}}.
\end{cousinweb}

\begin{grammarlens}[title={What you've built}]
One more word for telling what happened.
\end{grammarlens}

\subsection*{Your turn}
\begin{itemize}
\raggedright
  \item \emph{Say it:} \emph{ghoṣṇā}
  \item \emph{Say it:} \emph{ghoṣṇā}, once more
  \item \emph{Say it:} \emph{sarkār ne naī saṛak kī ghoṣṇā kī.} — The government announced a new road.
  \item \emph{Recall:} say the Hindi for a claim, then the Hindi for the public, then say \emph{ghoṣṇā} again
\end{itemize}

\begin{tcolorbox}[breakable,colback=teal!4,colframe=teal!35!black,title={Before you move on}]
What does \hi{घोषणा} mean? (\textquotedblleft{}A declaration; an official announcement\textquotedblright{}.) Say it once more.
\end{tcolorbox}
\section[\texorpdfstring{\hi{जानकारी} (jānkārī)}{जानकारी (jānkārī)}]{\hi{जानकारी} (jānkārī) — information}

\label{lesson:HI-C315-jankari}



\begin{quote}
Before the new one: say the Hindi for the public, then the Hindi for a declaration.
\end{quote}

\subsection*{You'll want to know: \hi{जानकारी}}
\textbf{\hi{जानकारी}} — \emph{jānkārī} — \textquotedblleft{}information\textquotedblright{}.

\begin{quote}
\hi{पुलिस} \hi{ने} \hi{पत्रकारों} \hi{को} \hi{कोई} \hi{जानकारी} \hi{नहीं} \hi{दी।}
\end{quote}

\begin{quote}
\emph{pulis ne patrakāroṁ ko koī jānkārī nahīṃ dī.}
\end{quote}

\begin{quote}
The police gave the journalists no information.
\end{quote}

Feminine.

\begin{cousinweb}
is \emph{jānkār}, one who knows, with the noun ending \emph{-ī}; \emph{jānkār} is built on \textbf{\hi{जानना}}, to know, from Sanskrit \emph{jānāti} (CDIAL 5193). \textbf{\hi{जानकारी}} is what a knower has: information.
\end{cousinweb}

\begin{grammarlens}[title={What you've built}]
One more word for telling what happened.
\end{grammarlens}

\subsection*{Your turn}
\begin{itemize}
\raggedright
  \item \emph{Say it:} \emph{jānkārī}
  \item \emph{Say it:} \emph{jānkārī}, once more
  \item \emph{Say it:} \emph{pulis ne patrakāroṁ ko koī jānkārī nahīṃ dī.} — The police gave the journalists no information.
  \item \emph{Recall:} say the Hindi for the public, then the Hindi for a declaration, then say \emph{jānkārī} again
\end{itemize}

\begin{tcolorbox}[breakable,colback=teal!4,colframe=teal!35!black,title={Before you move on}]
What does \hi{जानकारी} mean? (\textquotedblleft{}Information\textquotedblright{}.) Say it once more.
\end{tcolorbox}
\section[\texorpdfstring{\hi{टालना} (ṭālnā)}{टालना (ṭālnā)}]{\hi{टालना} (ṭālnā) — to put off; to dodge}

\label{lesson:HI-C315-ttaalna}



\begin{quote}
Before the new one: say the Hindi for a declaration, then the Hindi for information.
\end{quote}

\subsection*{You'll want to know: \hi{टालना}}
\textbf{\hi{टालना}} — \emph{ṭālnā} — \textquotedblleft{}to put off; to dodge\textquotedblright{}.

\begin{quote}
\hi{मंत्री} \hi{ने} \hi{पत्रकारों} \hi{का} \hi{सवाल} \hi{टाल} \hi{दिया।}
\end{quote}

\begin{quote}
\emph{mantrī ne patrakāroṁ kā savāl ṭāl diyā.}
\end{quote}

\begin{quote}
The minister dodged the journalists' question.
\end{quote}

Past, with \hi{ने} (he or she …): \hi{उसने} \hi{बात} \hi{टाल} \hi{दी}, \emph{usne bāt ṭāl dī}.

\begin{cousinweb}
\textbf{\hi{टालना}} (\emph{ṭālnā}) is the transitive partner of \textbf{\hi{टलना}}, from the same Sanskrit \emph{ṭalati} (CDIAL 5450): to make something move away from you, a question, a visit or a danger.
\end{cousinweb}

\begin{grammarlens}[title={What you've built}]
One more word for telling what happened.
\end{grammarlens}

\subsection*{Your turn}
\begin{itemize}
\raggedright
  \item \emph{Say it:} \emph{ṭālnā}
  \item \emph{Say it:} \emph{ṭālnā}, once more
  \item \emph{Say it:} \emph{mantrī ne patrakāroṁ kā savāl ṭāl diyā.} — The minister dodged the journalists' question.
  \item \emph{Recall:} say the Hindi for a declaration, then the Hindi for information, then say \emph{ṭālnā} again
\end{itemize}

\begin{tcolorbox}[breakable,colback=teal!4,colframe=teal!35!black,title={Before you move on}]
What does \hi{टालना} mean? (\textquotedblleft{}To put off; to dodge\textquotedblright{}.) Say it once more.
\end{tcolorbox}
